% =============================================================
%  Versão B — Analítica
%  Da Bobina ao Projeto: Reflexões sobre Processos de Fabricação
%  Materiais e Processos de Fabricação I — IFB Campus Samambaia
% =============================================================

\documentclass[12pt,a4paper,oneside,brazil,chapter=TITLE]{article}
\usepackage[T1]{fontenc}
\usepackage[utf8]{inputenc}
\usepackage[brazil]{babel}
\usepackage{times}
\renewcommand{\ttdefault}{cmtt}
\usepackage[a4paper,left=3cm,top=3cm,right=2cm,bottom=2cm,headheight=14pt,headsep=0.5cm]{geometry}
\usepackage{setspace}
\onehalfspacing
\usepackage{indentfirst}
\setlength{\parindent}{1.25cm}
\setlength{\parskip}{0pt}
\usepackage{titlesec}
\titleformat{\section}{\normalfont\normalsize\bfseries\MakeUppercase}{\thesection\quad}{0em}{}
\titleformat{\subsection}{\normalfont\normalsize\bfseries}{\thesubsection\quad}{0em}{}
\titleformat{\subsubsection}{\normalfont\normalsize\bfseries\itshape}{\thesubsubsection\quad}{0em}{}
\titlespacing*{\section}{0pt}{18pt}{12pt}
\titlespacing*{\subsection}{0pt}{18pt}{6pt}
\titlespacing*{\subsubsection}{0pt}{12pt}{6pt}
\newcommand{\secbreak}{\clearpage}
\let\oldsection\section
\renewcommand{\section}{\secbreak\oldsection}
\usepackage{fancyhdr}
\pagestyle{fancy}
\fancyhf{}
\fancyhead[R]{\fontsize{10}{12}\selectfont\thepage}
\renewcommand{\headrulewidth}{0pt}
\renewcommand{\footrulewidth}{0pt}
\fancypagestyle{plain}{\fancyhf{}\fancyhead[R]{\fontsize{10}{12}\selectfont\thepage}\renewcommand{\headrulewidth}{0pt}}
\usepackage{longtable}
\usepackage{booktabs}
\usepackage{array}
\usepackage{multirow}
\usepackage{tabularx}
\renewcommand{\arraystretch}{1.3}
\usepackage{caption}
\DeclareCaptionFormat{abnt}{\fontsize{10}{12}\selectfont#1#2#3}
\DeclareCaptionLabelSeparator{emdashsep}{ --- }
\captionsetup{format=abnt,justification=centering,singlelinecheck=false,labelsep=emdashsep,position=above,skip=3pt}
\newcommand{\fonte}[1]{\begin{flushleft}\fontsize{10}{12}\selectfont Fonte: #1\end{flushleft}}
\usepackage{enumitem}
\setlist{nosep, leftmargin=*}
\providecommand{\tightlist}{\setlength{\itemsep}{0pt}\setlength{\parskip}{0pt}}
\usepackage{tocloft}
\renewcommand{\contentsname}{SUMÁRIO}
\PassOptionsToPackage{hyphens}{url}
\usepackage[hidelinks,colorlinks=false,breaklinks=true]{hyperref}
\usepackage{quoting}
\quotingsetup{font=small,indentfirst=false,vskip=\baselineskip}
\usepackage{graphicx}
\usepackage{float}
\date{}

% =============================================================
\begin{document}

% ----- CAPA --------------------------------------------------
\begin{titlepage}
  \thispagestyle{empty}
  \begin{center}
    {\normalsize\bfseries
      Instituto Federal de Educação, Ciência e Tecnologia de Brasília --- IFB\\[0.3em]
      \textit{Campus} Samambaia\\[0.3em]
      Curso Superior de Tecnologia em Design de Produto
    }
    \vfill
    {\large\bfseries DA BOBINA AO PROJETO: REFLEXÕES SOBRE PROCESSOS DE FABRICAÇÃO DO AÇO A PARTIR DE VISITA TÉCNICA À GRAVIA INDÚSTRIA DE PERFILADOS DE AÇO}
    \vfill
    {\normalsize Brasília/DF\\2026}
  \end{center}
\end{titlepage}
\setcounter{page}{2}

% ----- FOLHA DE ROSTO ----------------------------------------
\begin{titlepage}
  \thispagestyle{empty}
  \begin{center}
    {\normalsize\bfseries VICTOR HUGO DA SILVA COSTA}
    \vfill
    {\large\bfseries DA BOBINA AO PROJETO: REFLEXÕES SOBRE PROCESSOS DE FABRICAÇÃO DO AÇO A PARTIR DE VISITA TÉCNICA À GRAVIA INDÚSTRIA DE PERFILADOS DE AÇO}
    \vfill
    \begin{flushright}
      \begin{minipage}{0.5\textwidth}
        \normalsize
        Relatório de visita técnica apresentado ao Curso Superior de Tecnologia em Design de Produto do Instituto Federal de Educação, Ciência e Tecnologia de Brasília --- IFB, \textit{Campus} Samambaia, como requisito parcial para a aprovação na disciplina de Materiais e Processos de Fabricação I.
        \vspace{1em}

        Professora: Prof.\textsuperscript{a} Keila Sanches
      \end{minipage}
    \end{flushright}
    \vfill
    {\normalsize Brasília/DF --- 2026}
  \end{center}
\end{titlepage}

% ----- RESUMO ------------------------------------------------
\clearpage
\thispagestyle{fancy}
\begin{center}
  {\normalsize\bfseries RESUMO}
\end{center}
\vspace{12pt}

\noindent
Este relatório registra e analisa a visita técnica realizada em 02 de junho de 2026 à Gravia Indústria de Perfilados de Aço, em Taguatinga, Distrito Federal, sob a perspectiva do designer de produto. O objetivo não é apenas descrever os processos observados, mas compreender o que eles implicam para as decisões de projeto: quais liberdades formais e dimensionais cada processo oferece, onde estão seus limites, e como a lógica industrial da produção em série condiciona --- e ao mesmo tempo viabiliza --- soluções projetuais. A partir da observação dos processos de corte a laser, corte a plasma, conformação por dobra e perfilação contínua, o relatório argumenta que compreender os processos de fabricação disponíveis é condição para projetar com realismo, especificidade e responsabilidade técnica.

\vspace{12pt}
\noindent
\textbf{Palavras-chave}: Design de produto; Processos de fabricação; Corte a laser; Conformação; Decisão de projeto.

% ----- SUMÁRIO -----------------------------------------------
\clearpage
\tableofcontents

% =============================================================
%  CORPO DO TEXTO
% =============================================================

\section{Introdução}

Existe uma lacuna habitual entre o desenho e a fábrica. O designer projeta formas; o fabricante enfrenta chapas, prensas e tolerâncias. Enquanto essa distância não é percorrida de forma concreta, os projetos tendem a oscilar entre dois extremos igualmente problemáticos: o excesso de restrição, em que o projetista se autocensura diante de processos que imagina mais limitantes do que são, e a ingenuidade formal, em que se especificam geometrias cuja viabilidade fabril nunca foi verificada.

A visita técnica à Gravia Indústria de Perfilados de Aço, realizada em 02 de junho de 2026, funcionou como uma travessia deliberada dessa lacuna. O objetivo não era observar os processos como espectador neutro, mas perceber o que eles comunicam ao designer: que decisões de projeto eles viabilizam, que limitações impõem, e como a lógica industrial da produção em série condiciona e ao mesmo tempo libera a criação de produtos em aço.

\section{A Lógica da Planta: Padronização, Série e Melhoria Contínua}

A primeira coisa que a Gravia comunica ao visitante é uma lógica: a planta não fabrica peças únicas por acidente, mas por exceção. O modelo produtivo é orientado à série, à repetição e à padronização. A organização em galpões especializados, cada um dedicado a uma família de processos, materializa fisicamente essa lógica. O fluxo do material pelo espaço é, em si, uma decisão de projeto industrial.

Essa orientação é reforçada pelas metodologias de gestão adotadas. O Lean Six Sigma combina a eliminação de desperdícios com o controle rigoroso da variabilidade, o que significa que cada processo é calibrado para produzir dentro de tolerâncias definidas, com o menor índice possível de desvio. A Gestão por Processos e Trabalho Padronizado (GPTW) complementa essa orientação ao documentar rotinas operacionais, reduzindo a dependência do conhecimento tácito individual.

Projetos que nascem compatíveis com a produção em série chegam ao mercado com custos menores, prazos mais curtos e qualidade mais consistente. Formas que respeitam os raios de dobra dos equipamentos, espessuras que se alinham ao portfólio de chapas disponíveis, geometrias que podem ser geradas por operações encadeadas de corte e conformação: tudo isso é mais do que otimização, é projeto com responsabilidade técnica. Projetar contra a lógica da fábrica é sempre possível; é apenas mais caro e mais lento.

\section{Da Bobina ao Perfil: o Fluxo como Decisão de Projeto}

O ponto de partida de tudo que é fabricado na Gravia é a bobina de aço. Esse dado aparentemente trivial carrega implicações importantes: a bobina define a largura máxima da chapa disponível para corte, e o aço em bobina é produzido em espessuras específicas, cujas escolhas condicionam quais processos subsequentes são aplicáveis e com quais parâmetros.

No Galpão 1, a bobina é desbobinada e cortada em chapas de 3 m ou 6 m, comprimentos que correspondem diretamente às faixas úteis das prensas dobradeiras disponíveis (Newton PSH-35030, com três metros, e Newton PSH-7030, com seis metros). Não é coincidência: o sistema de fabricação foi projetado de forma que as dimensões das chapas e as capacidades dos equipamentos sejam mutuamente compatíveis. Um componente concebido com comprimento próximo a 3 m ou 6 m será manufaturado com eficiência nessa planta; dimensões intermediárias implicarão sobras de material ou operações adicionais de ajuste.

Isso ficou evidente quando se percebeu que a maior parte das peças expostas no showroom respeitava, com folga mínima, justamente esses comprimentos. Reconhecer a compatibilidade entre projeto e processo como uma qualidade, e não como uma concessão, é parte do que se aprende ao observar o fluxo produtivo de dentro.

\section{Corte a Laser e Corte a Plasma: Precisão versus Produtividade}

A Gravia emprega dois processos térmicos de corte com características e aplicações distintas, e a comparação entre eles ajuda a entender decisões que o designer terá que tomar cedo.

O sistema de corte a laser (Newton LN 3015F) opera com elevada precisão dimensional e permite a programação direta de geometrias complexas via software. A geração do caminho de corte por arquivo CAD significa que o custo unitário de setup é praticamente nulo para novos formatos: desde que a geometria esteja dentro das capacidades do equipamento, qualquer forma pode ser produzida sem a fabricação de ferramental dedicado. Esse processo se traduz em liberdade geométrica com precisão, sendo o mais adequado para peças com recortes intrincados, furos de posicionamento ou formas que não podem ser obtidas por guilhotina ou puncionamento.

O corte a plasma (Oxipira Série Master, nos Galpões 4 a 6) opera com força diferente: sua vantagem está na capacidade de processar chapas de até 50 mm de espessura com produtividade elevada, ao custo de uma tolerância dimensional menor. Enquanto o laser trabalha melhor em espessuras reduzidas com exigência de precisão, o plasma é o processo adequado para componentes estruturais robustos, onde a variação dimensional de alguns milímetros é tolerável e a espessura do material exige uma fonte de energia que o laser não cobre com eficiência.

A escolha do processo de corte deve ser feita em função das características da peça (espessura, complexidade geométrica, tolerância requerida), não por hábito ou familiaridade. Um componente estrutural em chapa de 25 mm não pertence ao laser; uma peça de fixação com furo rosqueado de alta precisão não pertence ao plasma. Foi o que ficou mais claro na transição entre galpões.

\section{Conformação por Dobra: Liberdades e Limites Dimensionais}

As prensas dobradeiras hidráulicas constituem o processo central de conformação do aço na Gravia. Com dez equipamentos distribuídos entre três modelos no Galpão 1 (PSH-35030, PSH-7030 e PSH-25030) e mais um de grande porte nos galpões finais (PSH-55060, com capacidade para peças de até 6 m), a empresa cobre praticamente toda a faixa dimensional relevante para a construção civil.

O processo de dobra por prensa é, em essência, simples: a chapa é posicionada entre punção e matriz, e a força hidráulica produz a deflexão desejada. Mas essa simplicidade conceitual esconde uma série de parâmetros críticos. O raio mínimo de dobra depende da espessura e das propriedades mecânicas do material: aços mais resistentes exigem raios maiores para evitar trincas na zona de tração. O retorno elástico (\textit{springback}), a tendência do material de recuperar parcialmente a forma após a remoção da carga, deve ser compensado na programação do equipamento ou no projeto da peça. E os comprimentos de flange --- as dimensões mínimas das abas dobradas --- são determinados pelas dimensões dos ferramentais disponíveis na planta.

Há algo quase didático na prensa dobradeira que raramente aparece nos livros de fabricação: ela converte chapas planas em geometrias tridimensionais a um custo operacional surpreendentemente baixo, desde que o projeto seja compatível com sequências de dobra. Um projeto que decompõe volumes tridimensionais em séries de operações planares é um projeto que a fábrica absorve sem resistência. Perfis abertos em U, L, Z e C, obtidos com apenas uma operação por dobra, são a prova de que formas simples bem resolvidas têm uma lógica fabril elegante por trás.

\section{Perfilação Contínua: Quando a Padronização é o Produto}

O Galpão 3 opera sob uma lógica radicalmente diferente dos demais. Enquanto o Galpão 1 produz peças discretas, resultado de operações individuais sobre chapas, aqui o regime é contínuo: a bobina entra numa extremidade da linha e o perfil acabado sai pela outra, cortado automaticamente nos comprimentos especificados.

A perfiladeira observada opera a aproximadamente 32 m/min e processa espessuras entre 0,9 mm e 5 mm em larguras de até 1.200 mm. Os produtos resultantes --- trilhos e montantes para drywall, perfis metálicos e tubos de 2'' --- são componentes altamente padronizados, com seção transversal constante ao longo de todo o comprimento. Essa constância de seção é, simultaneamente, a característica que define o processo e o limite que precisa ser reconhecido.

A produção contínua é economicamente eficiente porque elimina o tempo morto entre operações: um rolo de 1.000 m é processado em menos de uma hora. Essa produtividade, porém, só se manifesta para perfis com seção constante. Qualquer variação ao longo do comprimento exigiria rolos de geometria variável, equipamento de custo e complexidade muito superiores. Perfis contínuos são componentes de altíssima competitividade de custo quando o projeto os comporta; variações de seção ao longo do comprimento precisarão ser obtidas por outros meios.

\section{Reflexões Finais: o que a Gravia Ensina ao Designer}

Sair de uma visita técnica com uma lista de modelos de equipamentos é levar de volta apenas a menor parte do que a experiência oferece.

O que ficou foi algo mais difícil de sistematizar: a percepção de que os processos de fabricação têm uma lógica própria que precede qualquer projeto, e que ignorá-la não é liberdade criativa, é apenas imprecisão. Cada processo tem seu custo, sua faixa de aplicação e seus limites. O corte a laser permite formas que nenhum ferramental mecânico poderia produzir com a mesma eficiência. A prensa dobradeira converte chapas planas em geometrias tridimensionais com custo operacional baixo, desde que o projeto seja compatível com sequências de dobra. A perfilação contínua viabiliza componentes padronizados a um custo impossível de replicar por outros meios. O corte a plasma atende ao que o laser não alcança em espessura.

Projetar bem em aço é trabalhar com essa gramática em mente, não para reproduzir o que a fábrica já faz, mas para explorar o que ela pode fazer quando o briefing é suficientemente informado.

% ----- REFERÊNCIAS -------------------------------------------
\section*{REFERÊNCIAS}
\addcontentsline{toc}{section}{REFERÊNCIAS}
\label{referencias}

\noindent
GRAVIA INDÚSTRIA DE PERFILADOS DE AÇO. \textbf{Visita técnica à unidade industrial de Taguatinga/DF}. [Observação direta]. Taguatinga, DF, 02 jun. 2026.

\vspace{6pt}
\noindent
IIDA, I.; GUIMARÃES, L. B. de M. \textbf{Ergonomia: Projeto e Produção}. 3. ed. São Paulo: Blucher, 2016.

\vspace{6pt}
\noindent
KALPAKJIAN, S.; SCHMID, S. R. \textbf{Manufacturing Engineering and Technology}. 7. ed. Upper Saddle River: Pearson, 2014.

\end{document}
